%! Characteristics of Rational Functions — Unit 5, Lesson 5.0 — Warm-Up
\documentclass[10pt]{article}
\usepackage{algebra2-article}
\usepackage{algebra2-boxes}
\newcommand{\TallMath}[1]{$\displaystyle #1\rule[-1.4em]{0pt}{3.2em}$}

\begin{document}

\pageheader{Unit 5, Lesson 5.0}{Warm-Up}
\namedateperiod

\vspace{0.10in}

\begin{notesbox}{Getting Ready: Skills We'll Use Today}
\small These three skills power today's lesson on reading rational graphs. Work quickly.

\vspace{4pt}
\textbf{1. Where is a denominator zero?} A fraction $\dfrac{P}{Q}$ has \emph{no value} when
$Q=0$. Set each denominator equal to zero and solve --- factor first when you can.
\begin{enumerate}[label=(\alph*), leftmargin=2.1em, itemsep=3pt]
  \item $x-1=0$ \quad$\Rightarrow$\quad $x=\blank{1.6cm}$
  \item $x^2-9=0$ \quad factored: \blank{3.0cm} \quad$\Rightarrow$\quad $x=\blank{2.4cm}$
  \item $x^2+x-6=0$ \quad factored: \blank{3.0cm} \quad$\Rightarrow$\quad $x=\blank{2.4cm}$
\end{enumerate}

\vspace{4pt}
\textbf{2. What happens \emph{near} a forbidden input --- and \emph{far} from it?} Let
$y=\dfrac{1}{x-1}$, which is undefined at $x=1$. Fill in both tables.
\begin{center}
\renewcommand{\arraystretch}{1.3}
\begin{tabular}{|c|c|c|c|c|}
\hline
\multicolumn{5}{|c|}{\small \emph{closing in on} $x=1$} \\ \hline
$x$ & $0.9$ & $0.99$ & $1.01$ & $1.1$ \\ \hline
$y$ & \blank{1.0cm} & \blank{1.0cm} & \blank{1.0cm} & \blank{1.0cm} \\ \hline
\end{tabular}
\hspace{0.25in}
\begin{tabular}{|c|c|c|c|}
\hline
\multicolumn{4}{|c|}{\small \emph{running off to the right}} \\ \hline
$x$ & $11$ & $101$ & $1001$ \\ \hline
$y$ & \blank{1.0cm} & \blank{1.0cm} & \blank{1.0cm} \\ \hline
\end{tabular}
\end{center}
As $x$ closes in on $1$, the outputs grow \blank{4.0cm}. \quad As $x$ runs off to the right, the
outputs close in on $y=\blank{1.0cm}$.

\vspace{4pt}
\textbf{3. Factor and cancel --- then check the \emph{original}.} Simplify
$\dfrac{x^2-4}{x-2}$.
\begin{enumerate}[label=(\alph*), leftmargin=2.1em, itemsep=3pt]
  \item Factor the numerator: $x^2-4=\blank{3.2cm}$, \; so after cancelling the expression is
        \blank{2.2cm}.
  \item Substitute $x=2$ into the \textbf{original} fraction. What do you get? \blank{4.4cm}
  \item So the simplified form and the original agree at every input \emph{except}
        $x=\blank{1.0cm}$.
\end{enumerate}

\end{notesbox}

\end{document}
